% Provenance (never printed): the markets-with-algebra strand of ../economics_series.md (grade 9: affine demand and
% supply; equilibrium by one equation; shifts as changes of intercept; the four directions; a real episode; the full
% model frame). Brief: sources/book1/BRIEFS.md, g9-01. Chart data: figdata/grade-9/01-demand-and-supply-equations
% (Eurostat prc_hicp_aind CP08202 / CP00 for France; WDI IT.CEL.SETS.P2 FRA; invented worked example and fish
% market). Ledger: g9:smith-1962. The brief's [A] fallback opening chart is used (France only: the one country with
% the mobile-telephone-equipment index from 2000); no photograph (the solar case is not used).
\chapter{Demand and Supply with Equations}\label{ch:g9:demand-and-supply-equations}

\begin{openingchart}{Mobile phones in France, 2000--2010: the real price of mobile telephones (their
price index divided by all consumer prices, 2000 = 100) against the number of mobile subscriptions per
100 people, one point a year.}{Eurostat, harmonised index of consumer \omterm{def:g7:money:price}{prices} (prc\_hicp\_aind),
mobile telephone equipment and all items; World \omterm{def:g8:banks-saving-interest:bank}{Bank}, World Development Indicators, IT.CEL.SETS.P2}
\begin{tikzpicture}
\begin{axis}[omrecord, xmin=40, xmax=100, ymin=0, ymax=110,
  xlabel={Mobile subscriptions per 100 people}, ylabel={Real price of a mobile phone, 2000 = 100},
  nodes near coords, point meta=explicit symbolic,
  visualization depends on={value \thisrow{pos} \as \labelpos},
  every node near coord/.append style={font=\footnotesize, anchor=\labelpos}]
\addplot[omDef, thick, mark=*, mark size=1.6pt] table[x=subs, y=price, meta=label]
  {figdata/grade-9/01-demand-and-supply-equations/out/phones.csv};
\end{axis}
\end{tikzpicture}
\end{openingchart}

% figdata: phones (2000: 100, 48.8; 2008: 13.9, 92.3; 2010: 9.3, 91.1)
In 2000, France counted $48.8$ mobile subscriptions per 100 people. In 2010 there were $91.1$. Over the
same years, the \omterm{def:g7:money:price}{price} of a mobile phone, compared with all other goods, fell to $9.3\,\%$ of its level
of 2000. The \omterm{def:g7:money:price}{price} went down and the number of subscriptions went up, year after year until 2008; after
that the \omterm{def:g7:money:price}{price} kept falling while subscriptions stayed near $92$ per 100 people.

Grade 7 explained such moves with demand and \omterm{def:g7:demand-and-supply-first-look:supply-curve}{supply curves} drawn from tables. With equations, the same
model gives exact numbers: the new \omterm{def:g7:money:price}{price}, the new quantity, and how the change is shared between buyers
and sellers.

\begin{recall}
At the \omterm{def:g7:demand-and-supply-first-look:equilibrium-price}{equilibrium price}, the \omterm{def:g7:demand-and-supply-first-look:quantity-demanded}{quantity demanded} equals the \omterm{def:g7:demand-and-supply-first-look:quantity-supplied}{quantity supplied}
(\cref{ch:g7:demand-and-supply-first-look}). An affine function has the form $y = m x + k$: its graph is
a straight line with slope $m$ and intercept $k$.
\end{recall}

\section{Demand as an affine function of the price}

\begin{definition}[Demand function]\label{def:g9:demand-and-supply-equations:demand-function}
A \emph{demand function}\index{demand function} gives the \omterm{def:g7:demand-and-supply-first-look:quantity-demanded}{quantity demanded} $Q^d$ for each \omterm{def:g7:money:price}{price} $p$.
In this chapter it is affine: $Q^d = a - b\,p$, with $a > 0$ and $b > 0$.
\end{definition}

The number $a$ is the \omterm{def:g7:demand-and-supply-first-look:quantity-demanded}{quantity demanded} at a \omterm{def:g7:money:price}{price} of zero; $b$ is how many fewer units buyers want for
each extra dalen of \omterm{def:g7:money:price}{price}. In the worked example of this chapter, buyers want $Q^d = 120 - 4p$ units a
day: 120 at a \omterm{def:g7:money:price}{price} of 0, 80 at a \omterm{def:g7:money:price}{price} of 10, none at a \omterm{def:g7:money:price}{price} of 30.

\section{Supply as an affine function}

\begin{definition}[Supply function]\label{def:g9:demand-and-supply-equations:supply-function}
A \emph{supply function}\index{supply function} gives the \omterm{def:g7:demand-and-supply-first-look:quantity-supplied}{quantity supplied} $Q^s$ for each \omterm{def:g7:money:price}{price} $p$.
In this chapter it is affine: $Q^s = c + d\,p$, with $d > 0$.
\end{definition}

In the worked example, sellers offer $Q^s = 20 + 6p$ units a day: 20 at a \omterm{def:g7:money:price}{price} of 0, and 6 more for
each extra dalen.

\begin{measure}[Reading a slope]
Economists draw the \omterm{def:g7:money:price}{price} on the vertical axis and the quantity on the horizontal axis. With $Q^d = 120
- 4p$, a rise of 1 dalen in the \omterm{def:g7:money:price}{price} lowers the quantity by 4 units: on the graph, the line goes 4
units left for each dalen up. A demand line with a small $b$ looks steep: buyers hardly change their
purchases when the \omterm{def:g7:money:price}{price} changes. A large $b$ gives a flat line.
\end{measure}

\section{Solving for the equilibrium}

At the \omterm{def:g7:demand-and-supply-first-look:equilibrium-price}{equilibrium price}, $Q^d = Q^s$. In the worked example:
\[
120 - 4p = 20 + 6p \quad\Longrightarrow\quad 100 = 10p \quad\Longrightarrow\quad p^* = 10,
\]
and the quantity is $Q^* = 120 - 4 \times 10 = 80$ (check: $20 + 6 \times 10 = 80$).

\begin{proposition}[Equilibrium of affine demand and supply]\label{prop:g9:demand-and-supply-equations:equilibrium}
If $Q^d = a - b\,p$ and $Q^s = c + d\,p$ with $b + d > 0$, the \omterm{def:g7:demand-and-supply-first-look:equilibrium-price}{equilibrium price} and quantity are
\[
p^* = \frac{a - c}{b + d}, \qquad Q^* = a - b\,p^*.
\]
\end{proposition}

\begin{proof}
At the equilibrium, $a - b\,p = c + d\,p$. Adding $b\,p - c$ to both sides gives $a - c = (b + d)\,p$;
dividing by $b + d$, which is not zero, gives $p^* = (a - c)/(b + d)$. The quantity is the \omterm{def:g7:demand-and-supply-first-look:quantity-demanded}{quantity
demanded} at that \omterm{def:g7:money:price}{price}.
\end{proof}

% figdata: facts (ex_p 10, ex_q 80, ex_new_p 8, ex_new_q 88)
\begin{omchart}{The worked example: demand $Q^d = 120 - 4p$, supply $Q^s = 20 + 6p$ (solid), and supply
after a fall in costs, $Q^s = 40 + 6p$ (dashed). The equilibrium moves from $p^* = 10$, $Q^* = 80$ to
$p^* = 8$, $Q^* = 88$.}{Invented}
\begin{tikzpicture}
\begin{axis}[omrecord, xmin=0, xmax=160, ymin=0, ymax=20, height=0.5\linewidth,
  xlabel={Quantity (units a day)}, ylabel={Price (dalens)}, xtick={0,40,80,88,120,160},
  xticklabel style={font=\scriptsize}]
\addplot[demand] table[x=qd, y=price] {figdata/grade-9/01-demand-and-supply-equations/out/example.csv};
\addplot[supply] table[x=qs, y=price, restrict x to domain=0:160] {figdata/grade-9/01-demand-and-supply-equations/out/example.csv};
\addplot[supply, shifted] table[x=qs_new, y=price, restrict x to domain=0:160] {figdata/grade-9/01-demand-and-supply-equations/out/example.csv};
\draw[guide] (axis cs:0,10) -- (axis cs:80,10) -- (axis cs:80,0);
\draw[guide] (axis cs:0,8) -- (axis cs:88,8) -- (axis cs:88,0);
\node[eqpoint] at (axis cs:80,10) {};
\node[eqpoint] at (axis cs:88,8) {};
\node[anchor=south west, font=\footnotesize, text=omDef] at (axis cs:122,0.3) {demand};
\node[anchor=south east, font=\footnotesize, text=omThm] at (axis cs:112,16.5) {supply};
\node[anchor=west, font=\footnotesize, text=omThm] at (axis cs:126,12) {new supply};
\end{axis}
\end{tikzpicture}
\end{omchart}

\section{Shifts as changes of intercept}

\begin{definition}[Shock]\label{def:g9:demand-and-supply-equations:shock}
A \emph{shock}\index{shock} is a change in something other than the \omterm{def:g7:money:price}{price} of the good that moves the
demand or the \omterm{def:g7:demand-and-supply-first-look:supply-curve}{supply curve}.
\end{definition}

A \omterm{def:g9:demand-and-supply-equations:shock}{shock} changes an intercept. A fall in costs that makes sellers offer 20 more units at every \omterm{def:g7:money:price}{price}
turns $Q^s = 20 + 6p$ into $Q^s = 40 + 6p$: the supply line moves right by 20 units. The new
equilibrium is $p^* = (120 - 40) \div 10 = 8$ and $Q^* = 120 - 4 \times 8 = 88$. The \omterm{def:g7:money:price}{price} falls by 2
dalens, the quantity rises by 8 units.

A \omterm{def:g9:demand-and-supply-equations:shock}{shock} to demand changes $a$ in the same way. Each of the four \omterm{def:g9:demand-and-supply-equations:shock}{shocks} moves the equilibrium in a
definite direction.

\begin{omfigure}
\begin{tikzpicture}[econ, x=0.55cm, y=0.55cm]
  \foreach \col/\row/\title/\kind in {0/0/Demand rises/dup, 9/0/Demand falls/ddown, 0/-9/Supply rises/sup, 9/-9/Supply falls/sdown} {
    \begin{scope}[shift={(\col,\row)}]
      \draw[econ axes] (0,5.5) node[above, font=\scriptsize] {price} |- (6.5,0) node[right, font=\scriptsize] {qty};
      \node[font=\small\bfseries, anchor=south] at (3.2,5.9) {\title};
      \draw[demand] (0.5,5) -- (5,0.5);
      \draw[supply] (0.5,0.5) -- (5,5);
    \end{scope}
  }
  % demand rises: new demand right
  \draw[demand, shifted] (1.5,5) -- (6,0.5);
  \node[eqpoint, fill=black!45] at (2.75,2.75) {}; \node[eqpoint] at (3.25,3.25) {};
  % demand falls
  \draw[demand, shifted] (9.0,4) -- (12.5,0.5);
  \node[eqpoint, fill=black!45] at (11.75,2.75) {}; \node[eqpoint] at (11,2) {};
  % supply rises
  \draw[supply, shifted] (1.5,-8.5) -- (6,-4);
  \node[eqpoint, fill=black!45] at (2.75,-6.25) {}; \node[eqpoint] at (3.25,-6.75) {};
  % supply falls
  \draw[supply, shifted] (9.5,-7.5) -- (13.0,-4);
  \node[eqpoint, fill=black!45] at (11.75,-6.25) {}; \node[eqpoint] at (11.25,-5.75) {};
  \node[align=left, font=\footnotesize, anchor=north west] at (0,-0.6) {price up, quantity up};
  \node[align=left, font=\footnotesize, anchor=north west] at (9,-0.6) {price down, quantity down};
  \node[align=left, font=\footnotesize, anchor=north west] at (0,-9.6) {price down, quantity up};
  \node[align=left, font=\footnotesize, anchor=north west] at (9,-9.6) {price up, quantity down};
\end{tikzpicture}

{\small The four \omterm{def:g9:demand-and-supply-equations:shock}{shocks}: the dashed line is the curve after the \omterm{def:g9:demand-and-supply-equations:shock}{shock}, the grey dot the old equilibrium and
the black dot the new one. A rise in demand raises both the
\omterm{def:g7:money:price}{price} and the quantity; a rise in supply lowers the \omterm{def:g7:money:price}{price} and raises the quantity.}
\end{omfigure}

\section{Four predictions and a test}

\begin{model}{Supply and demand}
\assumes Many buyers and many sellers of one good; the \omterm{def:g7:money:price}{price} is free to move; buyers buy less at a
higher \omterm{def:g7:money:price}{price} and sellers offer more.
\predicts A fall in costs (supply rises) lowers the \omterm{def:g7:money:price}{price} and raises the quantity; a rise in demand
raises both. With affine curves, the sizes follow from the slopes: a cost fall that adds $k$ units to
supply at every \omterm{def:g7:money:price}{price} lowers the \omterm{def:g7:money:price}{price} by $k/(b+d)$.
\tested The opening chart is consistent with a rise in supply: the real \omterm{def:g7:money:price}{price} of mobile phones fell
and the number of subscriptions rose. It is not a clean test, because the \omterm{def:g7:money:price}{price} is that of handsets
and the quantity that of subscriptions, two different goods; a test would need the \omterm{def:g7:money:price}{price} and the
number of handsets sold in one \omterm{def:g7:markets-and-prices:market}{market} (exercise 12 asks what else could produce the same picture). Laboratory \omterm{def:g7:markets-and-prices:market}{markets}, in which buyers and sellers are each told a private
value or cost and trade by calling out offers: \omterm{def:g7:money:price}{prices} end close to the crossing of the supply and
demand schedules (Smith, 1962).
\breaks Few sellers who can choose their \omterm{def:g7:money:price}{price} (\cref{ch:g8:competition-and-monopoly-first-look});
goods whose quality buyers cannot check; \omterm{def:g7:money:price}{prices} that are not allowed to move
(\cref{ch:g9:price-controls-taxes-subsidies}).
\end{model}

\begin{record}[Laboratory markets]
% ledger: g9:smith-1962
Classroom experiments published in 1962 (Smith, 1962) worked like this. Each buyer received a card
with the most they could pay for one unit; each seller a card with their cost. Nobody knew the others' cards,
so nobody knew where the supply and demand schedules crossed. Buyers and sellers called out bids and
offers until no more trades were possible; then the \omterm{def:g7:markets-and-prices:market}{market} reopened for a new ``trading day'' with the
same cards, usually four or five times. \omterm{def:g7:money:price}{Prices} and quantities ended close to the equilibrium of the
schedules, although no trader knew it. Many later experiments found the same, with as few as six to
eight traders.
\end{record}

\section{Exercises}

\begin{exercise}[$\star$]\label{exo:g9:demand-and-supply-equations:1}
Use the opening chart. Give the real \omterm{def:g7:money:price}{price} of a mobile phone and the number of subscriptions per 100
people in 2000 and in 2010.
\end{exercise}

\begin{exercise}[$\star$]\label{exo:g9:demand-and-supply-equations:2}
Demand is $Q^d = 200 - 5p$. Compute the \omterm{def:g7:demand-and-supply-first-look:quantity-demanded}{quantity demanded} at a \omterm{def:g7:money:price}{price} of 12 dalens.
\end{exercise}

\begin{exercise}[$\star$]\label{exo:g9:demand-and-supply-equations:3}
Demand is $Q^d = 90 - 3p$ and supply is $Q^s = 10 + 5p$. Compute the \omterm{def:g7:demand-and-supply-first-look:equilibrium-price}{equilibrium price} and quantity.
\end{exercise}

\begin{exercise}[$\star$]\label{exo:g9:demand-and-supply-equations:4}
A demand line passes through the points ($p = 0$, $Q = 300$) and ($p = 10$, $Q = 200$). Write its
equation.
\end{exercise}

\begin{exercise}[$\star\star$]\label{exo:g9:demand-and-supply-equations:5}
In the worked example, demand rises to $Q^d = 150 - 4p$; supply is still $Q^s = 20 + 6p$. Compute the
new equilibrium.
\end{exercise}

\begin{exercise}[$\star\star$]\label{exo:g9:demand-and-supply-equations:6}
In the worked example, supply falls to $Q^s = 6p$; demand is still $Q^d = 120 - 4p$. Compute the new
equilibrium.
\end{exercise}

\begin{exercise}[$\star\star$]\label{exo:g9:demand-and-supply-equations:7}
In the worked example, each unit becomes 2 dalens cheaper to make, so sellers now offer at a \omterm{def:g7:money:price}{price} $p$
what they offered before at $p + 2$: $Q^s = 20 + 6(p + 2)$. Compute the new \omterm{def:g7:demand-and-supply-first-look:equilibrium-price}{equilibrium price}. By how
much did the \omterm{def:g7:money:price}{price} fall: by the full 2 dalens?
\end{exercise}

\begin{exercise}[$\star\star$]\label{exo:g9:demand-and-supply-equations:8}
In the worked example, the \omterm{def:g7:money:price}{price} is held at 8 dalens. Compute the quantities demanded and supplied,
and the \omterm{def:g7:markets-and-prices:shortage}{shortage}.
\end{exercise}

\begin{exercise}[$\star\star$]\label{exo:g9:demand-and-supply-equations:9}
In a \omterm{def:g7:markets-and-prices:market}{market}, a new machine lowers costs and nothing else is said to change. A year later, both the
\omterm{def:g7:money:price}{price} and the quantity are higher. Does this fit the model? What else might have happened?
\end{exercise}

\begin{exercise}[$\star\star\star$]\label{exo:g9:demand-and-supply-equations:10}
In the worked example, demand rises to $Q^d = 150 - 4p$ and, at the same time, supply rises to $Q^s =
40 + 6p$. Compute the new equilibrium. Which of the \omterm{def:g7:money:price}{price} and the quantity could you predict without
the numbers?
\end{exercise}

\begin{exercise}[$\star\star\star$]\label{exo:g9:demand-and-supply-equations:11}
Supply is $Q^s = 20 + 6p$. A cost fall adds 12 units to supply at every \omterm{def:g7:money:price}{price}. Compute the fall in the
\omterm{def:g7:money:price}{price} when demand is $Q^d = 50 - p$, and when demand is $Q^d = 200 - 12p$. In which \omterm{def:g7:markets-and-prices:market}{market} do buyers get
most of the cost fall as a lower \omterm{def:g7:money:price}{price}?
\end{exercise}

\begin{exercise}[$\star\star\star$]\label{exo:g9:demand-and-supply-equations:12}
From 2008 to 2010, the opening chart shows the \omterm{def:g7:money:price}{price} still falling while subscriptions stopped rising.
What can the chart not say about which curves moved?
\end{exercise}

\section{Problem: The Brenmoor Fish Market}

\begin{problem}[{Weekend problem --- demand and supply lines from tables, the equilibrium, and a new
harbour}]
\label{pb:g9:demand-and-supply-equations:1}
% figdata: fish (qd = 1000 - 50p, qs = 50 + 75p, qs_new = 175 + 75p); facts (fish_p 7.6, fish_q 620;
% fish_new_p 6.6, fish_new_q 670)
Everything in this problem is invented. At the fish \omterm{def:g7:markets-and-prices:market}{market} of Brenmoor, buyers and fishermen would
trade these quantities each morning, in kilograms, at each \omterm{def:g7:money:price}{price} in dalens per kilogram:
\begin{center}
\begin{tabular}{lrrrrr}
\toprule
\omterm{def:g7:money:price}{Price} (dalens per kg) & 2 & 4 & 6 & 8 & 10 \\
\midrule
\omterm{def:g7:demand-and-supply-first-look:quantity-demanded}{Quantity demanded} (kg) & 900 & 800 & 700 & 600 & 500 \\
\omterm{def:g7:demand-and-supply-first-look:quantity-supplied}{Quantity supplied} (kg) & 200 & 350 & 500 & 650 & 800 \\
\bottomrule
\end{tabular}
\end{center}

\medskip
\noindent\textbf{Part I --- The lines.}
\begin{enumerate}
  \item By how many kilograms does the \omterm{def:g7:demand-and-supply-first-look:quantity-demanded}{quantity demanded} fall for each extra dalen of \omterm{def:g7:money:price}{price}?
  \item Write the \omterm{def:g9:demand-and-supply-equations:demand-function}{demand function} $Q^d = a - b\,p$.
  \item Write the \omterm{def:g9:demand-and-supply-equations:supply-function}{supply function} $Q^s = c + d\,p$.
  \item Compute the quantities demanded and supplied at 5 dalens per kilogram.
\end{enumerate}

\medskip
\noindent\textbf{Part II --- The equilibrium.}
\begin{enumerate}[resume]
  \item Solve for the \omterm{def:g7:demand-and-supply-first-look:equilibrium-price}{equilibrium price} and quantity.
  \item Check the quantity with the \omterm{def:g9:demand-and-supply-equations:supply-function}{supply function}.
  \item At 10 dalens per kilogram, compute the \omterm{def:g7:markets-and-prices:market-surplus}{market surplus}.
  \item At 6 dalens per kilogram, compute the \omterm{def:g7:markets-and-prices:shortage}{shortage}.
  \item Compute the fishermen's \omterm{def:g7:firms:revenue}{revenue} at the equilibrium.
\end{enumerate}

\medskip
\noindent\textbf{Part III --- A new harbour.} A new harbour lowers the fishermen's costs: at every \omterm{def:g7:money:price}{price}
they now bring 125 kg more, so $Q^s = 175 + 75p$.
\begin{enumerate}[resume]
  \item Compute the new \omterm{def:g7:demand-and-supply-first-look:quantity-supplied}{quantity supplied} at 6 dalens per kilogram.
  \item Solve for the new \omterm{def:g7:demand-and-supply-first-look:equilibrium-price}{equilibrium price} and quantity.
  \item By how much did the \omterm{def:g7:money:price}{price} and the quantity change?
  \item Compute the fishermen's new \omterm{def:g7:firms:revenue}{revenue}. Did it rise or fall?
\end{enumerate}
\end{problem}
